\chapter{JADWAL PENELITIAN}

Penelitian ini direncanakan berlangsung selama \textbf{4 bulan} dengan fokus pada perancangan, implementasi, pengujian, dan dokumentasi arsitektur perangkat lunak Anycademy. Jadwal disusun berdasarkan alur rekayasa perangkat lunak yang dijelaskan pada Bab III.

\section{Ringkasan Fase}

\begin{enumerate}
	\item \textbf{Fase 1 --- Perancangan Arsitektur (Bulan 1):} Analisis kebutuhan sistem, perancangan topologi arsitektur (\textit{modular monolith} + \textit{microservices}), perancangan skema basis data (PostgreSQL, MongoDB, Redis), dan perancangan API endpoints.

	\item \textbf{Fase 2 --- Implementasi (Bulan 1--3):} Pengembangan backend (Go/Gin), \textit{microservices} (Python FastAPI: IRT + AI), frontend (React/Vite), integrasi Redis/Asynq \textit{message queue}, implementasi \textit{multi-platform anti-cheat}, dan \textit{deployment} via Docker Compose.

	\item \textbf{Fase 3 --- Pengujian (Bulan 3--4):} \textit{Load} dan \textit{stress testing} dengan K6, pengujian efektivitas anti-kecurangan pada tiga platform, pengujian usabilitas antarmuka dengan \textit{System Usability Scale} (SUS), dan evaluasi performa \textit{pipeline} IRT.

	\item \textbf{Fase 4 --- Dokumentasi \& Pelaporan (Bulan 3--4):} Penulisan dokumentasi teknis, analisis hasil pengujian, penyusunan laporan tugas akhir, dan finalisasi naskah.
\end{enumerate}

\section{Gantt Chart}

\begin{table}[H]
	\centering
	\caption{Jadwal Penelitian Anycademy --- 4 Bulan}
	\begin{tabular}{|l|c|c|c|c|}
		\hline
		\multirow{2}{*}{\textbf{Kegiatan}}     & \multicolumn{4}{c|}{\textbf{Bulan}}                                        \\
		\cline{2-5}
		                                       & \textbf{1}                          & \textbf{2} & \textbf{3} & \textbf{4} \\
		\hline
		\multicolumn{5}{|c|}{\textbf{Fase 1: Perancangan}}                                                                  \\
		\hline
		Studi literatur \& analisis kebutuhan  & $\bullet$                           &            &            &            \\
		\hline
		Perancangan arsitektur sistem          & $\bullet$                           &            &            &            \\
		\hline
		Perancangan skema basis data           & $\bullet$                           &            &            &            \\
		\hline
		Perancangan API endpoints              & $\bullet$                           & $\bullet$  &            &            \\
		\hline
		\hline
		\multicolumn{5}{|c|}{\textbf{Fase 2: Implementasi}}                                                                 \\
		\hline
		Pengembangan backend (Go/Gin)          & $\bullet$                           & $\bullet$  & $\bullet$  &            \\
		\hline
		Pengembangan IRT Service (Python)      & $\bullet$                           & $\bullet$  &            &            \\
		\hline
		Pengembangan AI Service (Python)       & $\bullet$                           & $\bullet$  &            &            \\
		\hline
		Pengembangan frontend (React/Vite)     & $\bullet$                           & $\bullet$  & $\bullet$  &            \\
		\hline
		Integrasi Redis/Asynq queue            &                                     & $\bullet$  & $\bullet$  &            \\
		\hline
		Implementasi multi-platform anti-cheat &                                     & $\bullet$  & $\bullet$  &            \\
		\hline
		Deployment Docker Compose              &                                     & $\bullet$  & $\bullet$  &            \\
		\hline
		\hline
		\multicolumn{5}{|c|}{\textbf{Fase 3: Pengujian}}                                                                    \\
		\hline
		Load \& stress testing (K6)            &                                     &            & $\bullet$  & $\bullet$  \\
		\hline
		Pengujian efektivitas anti-cheat       &                                     &            & $\bullet$  & $\bullet$  \\
		\hline
		Pengujian usabilitas (SUS)             &                                     &            &            & $\bullet$  \\
		\hline
		Evaluasi performa IRT pipeline         &                                     &            & $\bullet$  &            \\
		\hline
		\hline
		\multicolumn{5}{|c|}{\textbf{Fase 4: Dokumentasi \& Pelaporan}}                                                     \\
		\hline
		Penulisan dokumentasi teknis           &                                     & $\bullet$  & $\bullet$  & $\bullet$  \\
		\hline
		Analisis hasil pengujian               &                                     &            & $\bullet$  & $\bullet$  \\
		\hline
		Penyusunan laporan TA                  &                                     &            & $\bullet$  & $\bullet$  \\
		\hline
		Finalisasi naskah                      &                                     &            &            & $\bullet$  \\
		\hline
	\end{tabular}
\end{table}

\noindent \textbf{Keterangan:} $\bullet$ = Kegiatan dilaksanakan pada bulan tersebut.

\section{Milestone Kritis}

\begin{itemize}
	\item \textbf{Akhir Bulan 1:} Desain arsitektur dan skema basis data selesai. Backend \textit{core} (auth, exam CRUD) berfungsi di lingkungan lokal.

	\item \textbf{Akhir Bulan 2:} Backend dan frontend terintegrasi penuh. IRT Service dan AI Service berfungsi. Sistem anti-cheat (\textit{KIOSK mode}, \textit{tamper detection}) terimplementasi pada web dan Android.

	\item \textbf{Akhir Bulan 3:} Seluruh fitur terimplementasi. Pengujian beban dengan K6 dan pengujian efektivitas anti-cheat selesai. Hasil pengujian terdokumentasi.

	\item \textbf{Akhir Bulan 4:} Laporan tugas akhir selesai, termasuk analisis hasil pengujian, kesimpulan, dan saran pengembangan. Naskah siap diserahkan.
\end{itemize}
